\section{Discussion}
\label{sec:discussion}

\subsection{Why retrieval gains transfer only in part}
\label{sec:why-transfer}
Eq.~\ref{eq:decomp} bounds the transfer ratio.
If P(EM$\mid$hit) and P(EM$\mid$miss) stayed at the values of MiniLM-L6-v2, $T$ would equal their difference, 0.33.
A hit is therefore worth a third of a correct answer for this reader, not a whole one.
The observed ratios, 0.19 to 0.25, fall below this value.
\emph{Observation:} the point estimates of P(EM$\mid$hit) for the stronger dense encoders (56.4 to 57.2) lie below that of MiniLM-L6-v2 (58.4), with overlapping intervals; we did not test this difference.
\emph{Hypothesis:} the extra hits of a stronger encoder fall on questions that the reader finds harder, so each added hit yields fewer correct answers.
Our design does not test this hypothesis.

Two further gaps limit transfer.
First, P(EM$\mid$hit) with five passages stays below Gold EM (65.3), which uses only the relevant passages.
\emph{Hypothesis:} non-relevant passages in the top five distract the reader, as \citet{leto2024toward} report for added non-gold documents.
Incomplete relevance labels are an alternative explanation: a ``miss'' may contain the answer in an unjudged passage.
Second, P(EM$\mid$miss) (24.4 to 26.0 for dense encoders) exceeds closed-book EM (14.9).
Unjudged answer-bearing passages, partial cues and matches on common short answers could each explain this; our design cannot separate them.

BM25 shows the highest P(EM$\mid$hit) (63.2).
\emph{Hypothesis:} BM25 retrieves the relevant passage mostly for questions with strong lexical overlap, which the reader also finds easier.
This is a selection effect, not evidence that BM25 passages are more useful to the reader.
An alternative is that the reader extracts answers more easily from lexically matched passages.
We cannot distinguish the two with our data.

\subsection{Why the largest encoders do not win}
\label{sec:why-large}
Within our setup, bge-large and Qwen3-Emb-0.6B do not retrieve better than e5-base-v2; they retrieve worse (\S\ref{sec:rq1}).
Their lack of an EM gain therefore follows from their retrieval, not from a failure of transfer.
Parameter count orders the BGE family but not the models across families.
\emph{Interpretation:} training data and objective matter more than size for this task, which agrees with the claim of \citet{baidya2026choosing}.
Our pipeline reproduces the model-card SciFact scores of the bge and e5 encoders (\S\ref{sec:models}).
In a post-hoc check, our SciFact nDCG@10 for Qwen3-Emb-0.6B (0.699) lies within 0.002 of its MTEB leaderboard value (0.697, earlier model revision).
A pipeline error is therefore an unlikely cause.
Prefixes and instructions are a possible cause that we did not vary (\S\ref{sec:stages}).
Transfer also compresses differences.
With $T$ near 0.25, a retrieval gap of 2 to 7\,pp corresponds to an expected EM gap of 0.5 to 1.8\,pp, below the minimum detectable difference of 3.4\,pp.
Differences in retrieval that are clearly significant can thus leave EM differences that 2,255 queries cannot resolve.

\subsection{Practical implication}
\label{sec:implication}
\emph{Interpretation:} for a fixed reader, a retrieval benchmark rank alone overstates the answer gain from a stronger encoder.
Within our setup, a Hit@5 gain of $x$\,pp came with an EM gain of $0.19x$ to $0.25x$\,pp relative to MiniLM-L6-v2, while encoding cost rose by up to 7.5$\times$ between e5-base-v2 and Qwen3-Emb-0.6B.
We suggest reporting the transfer ratio and measured cost alongside retrieval scores when an embedding model is chosen for RAG.
The ratio depends on the reader, the corpus and the question type, so it must be measured for each setup rather than taken from ours.

\subsection{Relation to prior work}
\label{sec:prior}
MTEB \citep{muennighoff2022mteb}, \citet{cirillo2026ge2} and \citet{baidya2026choosing} measure retrieval quality, and the latter two measure latency or cost.
Our retrieval results agree with their observation that size alone does not order models.
We add the downstream step: the retrieval ranking carries over to EM in order, but at 19 to 25\% of its size.
\citet{leto2024toward} find that lower search accuracy has minor effects on RAG performance; our ratios quantify a similar attenuation across seven retrieval methods rather than across search accuracy settings.
RAGAS \citep{es2023ragas} and ARES \citep{saadfalcon2024ares} evaluate answers but do not report how much of a retrieval gain reaches them; our design fixes the reader so that this quantity is identifiable.

\subsection{Threats to validity}
\label{sec:threats}
Our conclusions hold within one dataset, one reader, one corpus construction and one hardware setting.
Encoders use their model-card prefixes, so we do not compare prefix-free models.
The sub-corpus favours dense methods over BM25 and raises absolute scores.
EM is strict and undercounts correct paraphrases.
Training-data contamination of the encoders or the reader on NQ and Wikipedia is possible and not measurable here.
The judge is unvalidated, one annotator labels the validation sample, search is exact only, and the joined queries cover 65.3\% of the judged set.
The Limitations section lists each bound.
